\documentclass[14pt]{extarticle}

\usepackage[a4paper,margin=1in]{geometry}
\usepackage{amsmath,amssymb}
\usepackage{booktabs}
\usepackage{hyperref}

\def\mathbi#1{\textbf{\em #1}}

\title{Ridge Regression}
\author{}
\date{}

\begin{document}

\maketitle

Ridge regression is a supervised learning algorithm used for predicting a continuous target variable. It is an extension of linear regression that adds a regularization term to the loss function. This regularization term helps reduce overfitting and makes the model more stable when the feature matrix contains highly correlated features.

Suppose we are given a dataset consisting of $n$ observations and $p$ features. Let

\[
\mathbf{X} =
\begin{bmatrix}
x_{11} & x_{12} & \cdots & x_{1p} \\
x_{21} & x_{22} & \cdots & x_{2p} \\
\vdots & \vdots & \ddots & \vdots \\
x_{n1} & x_{n2} & \cdots & x_{np}
\end{bmatrix}
\]

be the feature matrix and

\[
\mathbi{y} =
\begin{bmatrix}
y_1 \\
y_2 \\
\vdots \\
y_n
\end{bmatrix}
\]

be the vector of target values.

As in linear regression, after absorbing the bias term into the feature matrix, the prediction function can be written as

\[
\hat{\mathbi{y}} = \mathbf{X}\mathbi{w},
\]

where $\mathbf{X} \in \mathbb{R}^{n \times p}$ is the feature matrix, $\mathbi{w} \in \mathbb{R}^{p}$ is the vector of model coefficients, and $\hat{\mathbi{y}} \in \mathbb{R}^{n}$ is the vector of predicted target values.

In ordinary linear regression, the loss function is

\[
L(\mathbi{w})
=
\frac{1}{2n}
\left\|
\mathbf{X}\mathbi{w}
-
\mathbi{y}
\right\|^2.
\]

Ridge regression adds an $L_2$ regularization term to penalize large model coefficients. Therefore, the ridge regression loss function is

\[
L(\mathbi{w})
=
\frac{1}{2n}
\left\|
\mathbf{X}\mathbi{w}
-
\mathbi{y}
\right\|^2
+
\frac{\lambda}{2}
\left\|
\mathbi{w}
\right\|^2,
\]

where $\lambda \geq 0$ is the regularization parameter. A larger value of $\lambda$ applies a stronger penalty to large coefficients.

Equivalently, the loss function can be written as

\[
L(\mathbi{w})
=
\frac{1}{2n}
(\mathbf{X}\mathbi{w}-\mathbi{y})^\top
(\mathbf{X}\mathbi{w}-\mathbi{y})
+
\frac{\lambda}{2}
\mathbi{w}^\top \mathbi{w}.
\]

Taking the derivative with respect to $\mathbi{w}$, we obtain

\begin{align*}
\frac{\partial L(\mathbi{w})}{\partial \mathbi{w}}
&=
\frac{\partial}{\partial \mathbi{w}}
\left[
\frac{1}{2n}
(\mathbf{X}\mathbi{w}-\mathbi{y})^\top
(\mathbf{X}\mathbi{w}-\mathbi{y})
+
\frac{\lambda}{2}
\mathbi{w}^\top \mathbi{w}
\right] \\
&=
\frac{1}{n}
\mathbf{X}^\top
(\mathbf{X}\mathbi{w}-\mathbi{y})
+
\lambda \mathbi{w}.
\end{align*}

Setting the derivative equal to $\mathbf{0}$ gives

\begin{align*}
\frac{1}{n}
\mathbf{X}^\top
(\mathbf{X}\mathbi{w}-\mathbi{y})
+
\lambda \mathbi{w}
&=
\mathbf{0} \\
\Leftrightarrow \mathbf{X}^\top
(\mathbf{X}\mathbi{w}-\mathbi{y})
+
n\lambda \mathbi{w}
&=
\mathbf{0} \\
\Leftrightarrow \mathbf{X}^\top\mathbf{X}\mathbi{w}
-
\mathbf{X}^\top\mathbi{y}
+
n\lambda \mathbi{w}
&=
\mathbf{0} \\
\Leftrightarrow \left(
\mathbf{X}^\top\mathbf{X}
+
n\lambda \mathbf{I}
\right)
\mathbi{w}
&=
\mathbf{X}^\top\mathbi{y}.
\end{align*}

Assuming $\mathbf{X}^\top\mathbf{X} + n\lambda \mathbf{I}$ is invertible, we obtain

\[
\mathbi{w}
=
\left(
\mathbf{X}^\top\mathbf{X}
+
n\lambda \mathbf{I}
\right)^{-1}
\mathbf{X}^\top\mathbi{y}.
\]

This is the closed-form solution for ridge regression.

When $\lambda = 0$, ridge regression becomes ordinary linear regression. When $\lambda$ increases, the model coefficients are pushed closer to $0$, which can reduce overfitting.

\end{document}